\documentclass[10pt,a4paper]{amsart}
\usepackage{amsmath,amssymb,amsthm,graphicx}
\usepackage{verbatim}
\newenvironment{reference}{\comment}{\endcomment}
%\newenvironment{reference}{}{}
\newenvironment{slogan}{\comment}{\endcomment}
\newenvironment{history}{\comment}{\endcomment}

% For commutative diagrams we use Xy-pic
\usepackage[all]{xy}

% We use 2cell for 2-commutative diagrams.
\xyoption{2cell}
% Theorem environments.
%
\theoremstyle{plain}
\newtheorem{theorem}[subsection]{Theorem}
\newtheorem{proposition}[subsection]{Proposition}
\newtheorem{lemma}[subsection]{Lemma}

\theoremstyle{definition}
\newtheorem{definition}[subsection]{Definition}
\newtheorem{example}[subsection]{Example}
\newtheorem{exercise}[subsection]{Exercise}
\newtheorem{situation}[subsection]{Situation}

\theoremstyle{remark}
\newtheorem{remark}[subsection]{Remark}
\newtheorem{remarks}[subsection]{Remarks}

\numberwithin{equation}{subsection}

% Macros
%
\def\lim{\mathop{\mathrm{lim}}\nolimits}
\def\colim{\mathop{\mathrm{colim}}\nolimits}
\def\Spec{\mathop{\mathrm{Spec}}}
\def\Hom{\mathop{\mathrm{Hom}}\nolimits}
\def\Ext{\mathop{\mathrm{Ext}}\nolimits}
\def\SheafHom{\mathop{\mathcal{H}\!\mathit{om}}\nolimits}
\def\SheafExt{\mathop{\mathcal{E}\!\mathit{xt}}\nolimits}
\def\Sch{\mathit{Sch}}
\def\Mor{\mathop{\mathrm{Mor}}\nolimits}
\def\Ob{\mathop{\mathrm{Ob}}\nolimits}
\def\Sh{\mathop{\mathit{Sh}}\nolimits}
\def\NL{\mathop{N\!L}\nolimits}
\def\CH{\mathop{\mathrm{CH}}\nolimits}
\def\proetale{{pro\text{-}\acute{e}tale}}
\def\etale{{\acute{e}tale}}
\def\QCoh{\mathit{QCoh}}
\def\Ker{\mathop{\mathrm{Ker}}}
\def\Im{\mathop{\mathrm{Im}}}
\def\Coker{\mathop{\mathrm{Coker}}}
\def\Coim{\mathop{\mathrm{Coim}}}

% Boxtimes
%
\DeclareMathSymbol{\boxtimes}{\mathbin}{AMSa}{"02}

%
% Macros for moduli stacks/spaces
%
\def\QCohstack{\mathcal{QC}\!\mathit{oh}}
\def\Cohstack{\mathcal{C}\!\mathit{oh}}
\def\Spacesstack{\mathcal{S}\!\mathit{paces}}
\def\Quotfunctor{\mathrm{Quot}}
\def\Hilbfunctor{\mathrm{Hilb}}
\def\Curvesstack{\mathcal{C}\!\mathit{urves}}
\def\Polarizedstack{\mathcal{P}\!\mathit{olarized}}
\def\Complexesstack{\mathcal{C}\!\mathit{omplexes}}
% \Pic is the operator that assigns to X its picard group, usage \Pic(X)
% \Picardstack_{X/B} denotes the Picard stack of X over B
% \Picardfunctor_{X/B} denotes the Picard functor of X over B
\def\Pic{\mathop{\mathrm{Pic}}\nolimits}
\def\Picardstack{\mathcal{P}\!\mathit{ic}}
\def\Picardfunctor{\mathrm{Pic}}
\def\Deformationcategory{\mathcal{D}\!\mathit{ef}}
\usepackage{cleveref}
\crefname{originallemma}{original lemma}{original lemmas}
\Crefname{originallemma}{Original lemma}{Original lemmas}
\makeatletter
\newlabel{lemma-double-complex-gives-resolution}{{0133}{source}{}{sourcepointer-1}{}}
\newlabel{lemma-double-complex-gives-resolution@cref}{{[originallemma][0][]0133}{[1][0][]source}{}{}{}}
\makeatother
\title{The total complex of a resolution computes the cohomology of the resolved complex}
\author{The Stacks Project Authors}
\date{stacks-project:homology.tex:local-snapshot-20261012; GNU FDL 1.2 or later}
\begin{document}
\maketitle
\noindent\textit{Editorial note.} All delivered mathematical text is the author's own, copied from the pinned Stacks Project chapter file homology.tex; no statement or proof step was rewritten, reordered or removed.\par\medskip
\noindent\textit{Editorial note.} The proof is the author's complete argument; the one sentence that points at the companion lemma is a reduction to a named same-chapter statement carried on this card as a navigation stub, and the general case is then handled by the author's own truncation argument. The closing sentence 'is injective and similar arguments' refers back to the argument displayed immediately before it and is preserved verbatim.\par\medskip
\noindent\textit{Editorial note.} Every labelled reference occurring inside the counted spans is carried in assembly.external\_references as a non-load-bearing navigation pointer (same-chapter ones at their definition site, cross-chapter ones at the invocation inside the counted span), so no load-bearing mathematics is replaced by a stub. No bibliography citation occurs inside the counted spans.\par\medskip
\noindent\textit{Editorial note.} The same-name check for this card scanned every staged Stacks chapter of this repository -- all sources/stacks-*/pinned-original-source.tex files, the card's own chapter excluded -- and the labels found are recorded in dedup.same\_name\_cross\_chapter on the card.\par\medskip
\noindent\textit{Editorial note.} The preamble is the chapter's own native preamble excerpt already used for this source: the comment-type reference, slogan and history environments (native preamble lines 8-13), the Xy-pic setup (native preamble lines 14-18) and the native theorem declarations with the mathematical macros (native preamble lines 154-217, the actual last line of the file).\par\medskip
\noindent\textit{Allowed standard prerequisites (editorial).}\begin{itemize}
\item The bounded below case. Same chapter, Lemma lemma-double-complex-gives-resolution. Supplies the assertion in the case where the resolution is concentrated in degrees bounded below, to which the general case is reduced by truncation.
\item Stupid truncations of a complex. Same chapter, the definition of the stupid truncation of a complex in degrees at least t. Supplies the truncated complexes whose total complexes carry the cohomology classes used in the reduction.
\item The total complex of a double complex. Same chapter, the construction of the total complex. Supplies the object Tot(A) receiving the map from M and the cohomology classes considered in the proof.
\end{itemize}\medskip
\ifdefined\hypertarget\hypertarget{sourcepointer-1}{}\fi
\noindent\textit{Source reference.} Original lemma 0133: sources/stacks-homology/pinned-original-source.tex, source lines 6916--6916. Navigation pointer invoked at proof line 7068; the counted span contains the pointer only and the cited result is not reproduced inside it, so the pointer is registered as a non-load-bearing navigation stub in the established form for Stacks material. No bibliography citation is involved.\par\medskip
\begin{lemma}
\label{lemma-right-resolution-gives-qis}
Let $M^\bullet$ be a complex of abelian groups. Let
$$
0 \to M^\bullet \to A_0^\bullet \to A_1^\bullet \to A_2^\bullet \to \ldots
$$
be an exact complex of complexes of abelian groups. Set
$A^{p, q} = A_p^q$ to obtain a double complex.
Then the map $M^\bullet \to \text{Tot}(A^{\bullet, \bullet})$
induced by $M^\bullet \to A_0^\bullet$ is a quasi-isomorphism.
\end{lemma}

\begin{proof}
If there exists a $t \in \mathbf{Z}$ such that $A_0^q = 0$ for $q < t$, then
this follows immediately from
Lemma \ref{lemma-double-complex-gives-resolution}
(with $p$ and $q$ swapped as in the final statement of that lemma).
OK, but for every $t \in \mathbf{Z}$ we have a complex
$$
0 \to
\sigma_{\geq t}M^\bullet \to
\sigma_{\geq t}A_0^\bullet \to
\sigma_{\geq t}A_1^\bullet \to
\sigma_{\geq t}A_2^\bullet \to \ldots
$$
of stupid truncations. Denote $A(t)^{\bullet, \bullet}$ the corresponding
double complex. Every element $\xi$ of $H^n(\text{Tot}(A^{\bullet, \bullet}))$
is the image of an element of $H^n(\text{Tot}(A(t)^{\bullet, \bullet}))$
for some $t$ (look at explicit representatives of cohomology classes).
Hence $\xi$ is in the image of $H^n(\sigma_{\geq t}M^\bullet)$.
Thus the map $H^n(M^\bullet) \to H^n(\text{Tot}(A^{\bullet, \bullet}))$
is surjective. It is injective because for all $t$ the map
$H^n(\sigma_{\geq t}M^\bullet) \to H^n(\text{Tot}(A(t)^{\bullet, \bullet}))$
is injective and similar arguments.
\end{proof}

\end{document}
